% Author revision: insert as an appendix or a subsection of the construction.
% Requires amsmath, amssymb, booktabs, array and hyperref (already in manuscript).
\subsection{Specification of the retained V4.6 mechanism and its revenue}
\label{sec:active-reference-specification}

The retained V4.6 lower-bound mechanism is $M_J$, including the final joint price
release. Its reference before the new modifications is $M_{4.5}$. This
specification fixes their parameters, order of composition and exceptional
reports. All source paths below are relative to
\path{certificate/joint_residual_screening_lower_bound/source/} in the
archive. A profile is ordered as
$\omega=(w_1,w_2,v_1,v_2)\in[0,1]^4$: $w$ is bidder 1's report and $v$ is
bidder 2's report, each in physical item order. A menu option $(x;p)$ gives
expected allocation $x\in[0,1]^2$, payment $p$, and utility $w\cdot x-p$
or $v\cdot x-p$ for the respective bidder.
Masks $0,1,2,3$ mean empty, item 1, item 2 and both items.

\paragraph{Constants and retained tariffs.}
Put
\[
\begin{gathered}
 a=\frac{159}{250},\quad b=\frac{91}{100},\quad c=b-a=\frac{137}{500},\\
 s_0=\frac{1137}{1000},\quad s=\frac{142}{125},\quad
 d=s-a=\frac12,\quad q=s-b=\frac{113}{500}.
\end{gathered}
\]
The earlier deterministic construction uses $s_0$, hence
$d_0=501/1000$ and $q_0=227/1000$. This historical mechanism changes the shared
allocation cost to $s=s_0-1/1000$ and retains the V3 price increments as
functions of the opponent report. In particular, the constants in those
frozen functions are not replaced by $d,q$.
For an opponent report $u$, define
\[
\begin{split}
 H(u)&=\max\{0,u_1-a,u_2-a,u_1+u_2-b\},\\
 m^r(u)&=(0,H(u)+\min\{a,r-u_2\},
                 H(u)+\min\{a,r-u_1\},H(u)+b).
\end{split}
\]
Let $P^0(u)$ be the four prices returned by
\path{V3/verifier/joined_threshold.py:menu(u)}, using
\path{V3/verifier/baseline_mechanism.py:menu(u)['prices']} when the first
function returns \texttt{None}. These are explicit piecewise algebraic
functions: the joined rule applies on
$d_0\le t\le1$, $0\le\rho\le b-\min(t,a)$, where
$t=\max(u)$, $\rho=\min(u)$; its successive branches use the tests
$t\le a$, $(t-q_0)^2\le r_0$, $(t-q_0)^2<r_1$, and $t\le2/3$, with
$r_0=b^2-(2/3+c^2)$ and $r_1=4c/3-2/9$ exactly as in that function.
The remaining affine branch is
\path{V3/verifier/stationary_s.py:prices_for_t(t)}.
The finite tariff data are
\path{V3/certificate/baseline_mechanism.json}: 20 common rows, 41 bundle
rows and 8 item rows, in their stored order. The literal first-match and
endpoint rules are \path{baseline_mechanism.py:fee_rows} and
\path{baseline_mechanism.py:menu} in the same verifier directory.
Thus the retained prices used by both bidders are exactly
\begin{equation}\label{eq:active-retained-prices}
 P^s(u)=m^s(u)+P^0(u)-m^{s_0}(u).
\end{equation}
This is \path{V4_02/verifier/split_cost_candidate.py:retained_menu}
at \texttt{CANDIDATE\_THETA}$=-1/1000$.

\paragraph{Shared selection and the reference $M_{4.5}$.}
First maximize reported total value minus cost over the ordered mask pairs
\[
 (0,0),(1,0),(2,0),(3,0),(0,1),(0,2),(0,3),(1,2),(2,1),
\]
with respective costs $0,a,a,b,a,a,b,s,s$, choosing the first maximizer.
This is the retained source's literal \texttt{base\_outcome\_order}, including
$(3,0)$ before $(0,1)$; it fixes the tie convention for this construction.
For each bidder maximize utility at $P^s(u)$: choose empty if the maximum
is zero; otherwise retain that bidder's shared mask if it maximizes utility;
otherwise choose the smallest maximizing mask contained in the shared mask.
This selection rule is part of the mechanism.

Next override bidder 2's menu when
$w\in Q=\{\max(w)\le A,\ w_1+w_2\le b\}$, where $A=2/3$ and
$b_0=(4-\sqrt2)/3$. If $t=\max(w)\le d$, its singleton and bundle prices
are $(A,A,b_0)$, with mask priority $(0,1,2,3)$. If $t>d$, put
$k=t-q$, $C_0=5/6+3k^2/4$ and use aligned prices
$(A,C_0-k,C_0)$ for the scarce singleton, safe singleton and bundle,
with priority (empty, safe, scarce, bundle). The scarce item is the item
at the opponent's larger coordinate. The source predicate is
\texttt{u[0] > u[1]} for scarce item 1 and item 2 otherwise; equal
coordinates cannot occur in a constrained $Q$ row because $t>d$ and
$\sum u\le b<2d$. These closed-$Q$ rules are
\path{split_cost_candidate.py:screening_menu}; the resulting evaluator is
\path{V4_02/verifier/split_cost_candidate.py:candidate}.

Finally, for each bidder whose opponent satisfies
$A<t<T=613/750$ and $u_1+u_2\le b$, replace its menu by
\begin{equation}\label{eq:active-lottery-menu}
 (0,0;0),\ (0,1;d+\delta),\ (1,0;t),\ (1,1;t+c+\delta),\
 (1,\beta;t+\beta c),
\end{equation}
in the displayed priority, where
\[
 \delta=q+\frac{3q^2}{4}-\frac{3qt}{2}+\frac{(3t-2)^2}{16},
 \qquad \beta=\frac{3t-2}{3t-2+2\delta}.
\]
The first aligned coordinate is the scarce item; reverse allocation
coordinates when $u_1<u_2$. This gives $M_{4.5}$, evaluated by
\path{V4_5/verifier/residual_lottery.py:mechanism}.
The endpoints $t=A,T$ retain the preceding rule; the face $\sum u=b$
is included in the lottery rule.

\paragraph{Ordered modifications defining $M_J$.}
Apply the following operations to $M_{4.5}$, retaining every earlier row
unless it is explicitly replaced. First, for both bidders extend
\eqref{eq:active-lottery-menu} to $A<t<T$, $b-t<\rho<c$; retain the
preceding rule on $\rho=c$. Second, for bidder 1's opponents in
$F=\{\max(v)\le d,\ \sum v\le b_0\}$ use prices $(A,A,b_0)$.
For opponents in $G=\{\max(v)\le d,\ b_0<\sum v<b\}$ use
$(a,a,\sum v)$. Both use mask priority $(0,1,2,3)$.
Third, on bidder 2's $Q$ rows replace the bundle price by
$\max(b_0,\sum w)$ if $\max(w)\le d$.
On a constrained row, if $\sum w>C_0$, replace the aligned prices by
$(A,\sum w-k,\sum w)$; otherwise retain its preceding $Q$ menu.
Keep the $Q$ priorities above. The boundaries $\sum v=b_0,b$ retain their
preceding bidder-1 rules, and equality $\sum w=C_0$ retains the earlier
constrained row. These operations, including the intermediate bundle-only
response before the coupled response, are the successive
\texttt{mechanism} functions in
\path{V4_6/verifier/outer_lottery_strip.py},
\path{free_bidder_one.py}, \path{outer_bundle_exchange.py} and
\path{outer_bundle_reoptimized.py}, all in that verifier directory.
Call the resulting mechanism $M_6$.

Set $\varepsilon=9/10000$ and use the closed rectangles
\[
 W=[7/10,71/100]\times[0,1/5],\qquad
 S=[7/10-4\varepsilon,71/100]\times[c-2\varepsilon,c+4\varepsilon].
\]
When $v\in S$, charge bidder 1 an additional $\varepsilon$ on every
nonempty option of its $M_6$ menu: choose empty when its previous maximum
utility is at most $\varepsilon$, and otherwise retain its selected option.
When $w\in W$, offer bidder 2 the five options in
\eqref{eq:active-lottery-menu}, lowering only the lottery payment by
$\varepsilon$, with the same first-maximizer priority. This last operation
uses the displayed physical orientation only. The final evaluator is
\path{V4_6/verifier/price_joint_reallocation.py:mechanism(profile)}.
Its \texttt{allocations} and \texttt{payments} fields define $M_J$.
The real mechanism interprets every displayed and source-defined branch
over the reals; the executable evaluates rational reports using exact
rational and quadratic comparisons. No floating-point extension or numerical
tie tolerance defines the mechanism.

\paragraph{Exact revenue.}
For each specified evaluator $E$, let $p_i^E(w,v)$ be its selected payment.
With four independent uniform values, define the exact number
\begin{equation}\label{eq:active-payment-integral}
 \mathcal R(E)=\int_0^1\!\int_0^1\!\int_0^1\!\int_0^1
 [p_1^E(w_1,w_2,v_1,v_2)+p_2^E(w_1,w_2,v_1,v_2)]
 \,dv_2\,dv_1\,dw_2\,dw_1.
\end{equation}
In particular $R_{4.5}=\mathcal R(M_{4.5})$ and
$R_J=\mathcal R(M_J)$. The following table fixes the reduced-integral
entrypoints; it does not define revenue through a saved decimal.

\begingroup\small
\renewcommand{\arraystretch}{1.18}
\begin{center}
\begin{tabular}{@{}>{\raggedright\arraybackslash}p{0.18\linewidth}>{\raggedright\arraybackslash}p{0.76\linewidth}@{}}
\toprule
Quantity & Exact definition and reduction in the source directory\\
\midrule
Historical base $R_{02}$ &
$\mathcal R(\texttt{split\_cost\_candidate.candidate})$.
\path{V4_02/verifier/split_cost_revenue.py:calculate} reduces its change
from V3.1 using \texttt{BASE\_D}, \texttt{base\_Q}, \texttt{inner\_Q}
and \texttt{moments}; V3.1 is
\path{V3_1/verifier/constrained_candidate.py:mechanism}.\\
$R_{4.5}-R_{02}$ &
$4\int_A^{t_o}(b-t)g_1(t)\,dt+
 4\int_{t_o}^{T}(b-t)g_2(t)\,dt$;
\path{V4_5/verifier/independent_lottery.py:reconstruct}.\\
$R_6-R_{4.5}$ &
$\int_{[0,1]^4}\sum_i(p_i^{M_6}-p_i^{M_{4.5}})$;
\path{V4_6/verifier/independent_revenue.py:calculate}
returns \texttt{refined\_addition\_coefficients} in the basis
$(1,\sqrt2,\sqrt{23})$.\\
$J=R_J-R_6$ &
$\int_{[0,1]^4}\sum_i(p_i^{M_J}-p_i^{M_6})$;
\path{V4_6/verifier/price_joint_revenue.py:calculate}
reduces this integral to a rational constant and three logarithms.\\
\bottomrule
\end{tabular}
\end{center}
\endgroup
Here $\delta_o(t)=1/2-c+3q_0^2/4-3q_0t/2$,
$t_o=(1/2-c+3q_0^2/4)/(3q_0/2)$,
$g_1=(\delta-\delta_o)^2+\delta_o(3t-2)^2/8$ and $g_2=\delta^2$.
The V4.02 reduction uses its explicitly integrated polynomial on
$[s_0,s]$ and the fixed tariff moments; its historical base ultimately
shares the earlier affine-polytope revenue calculation.
\path{V4_6_archive_audit/revenue/fresh_exact_gap.py:algebraic_increment}
and \texttt{joint\_increment} independently reconstruct the two new
V4.6 additions. They retain the V4.5 base enclosure and therefore do not
constitute a second integration of the entire historical base.
The rational radical and logarithm bounds in that file and
\path{V4_6/verifier/consolidated_bounds.py:verify} give
\[
 0.8758198541484224553460<R_J<0.8758198541484224553461.
\]
This definition and enclosure concern the constructed mechanism. They do
not assert unrestricted optimality or a matching full-auction certificate.
